\subsection{在实值函数的延拓中的应用}

由于 \(\mathbf{R}\) 可以看作 \(\tilde{\mathbf{R}}\) 的一个子集，每个集 \(\mathbf{E}\) 到 \(\mathbf{R}\) 中的映射（即 \(\mathbf{E}\) 上的每个实值函数）都可以看作 \(\mathbf{E}\) 到 \(\tilde{\mathbf{R}}\) 中的映射；特别地，若 \(\mathbf{E}\) 是拓扑空间 \(\mathbf{F}\) 的子集，且 \(f\) 是 \(\mathbf{E}\) 到 \(\mathbf{R}\) 中的映射，则可能发生这样的情形：在闭包 \(\overline{\mathbf{E}}\) ——\(\mathbf{E}\) 的闭包——的某些点处，\(f(x)\) 趋于极限 \(\infty\) （当 \(x\) 保持属于 \(\mathbf{E}\) 而趋于这些点之一时）；这时我们就可以把函数 \(f\) 连续延拓，在这些点处给它值 \(\infty\)（第 I 章，§ 8，第 5 小节，定理 1）。

特别地考察 E 是 \(R^{n}\) 的子集的情形，这时空间 \(R^{n}\) 本身看作嵌入在射影空间 \(P_{n}\) 中；若假设无穷远超平面是 \(x_{0}=0\) ，定义在 E 上的实值函数 f 就可以等同于映射

\[
(x _ {0}, x _ {1}, x _ {2}, \ldots , x _ {n}) \rightarrow f \left(\frac {x _ {1}}{x _ {0}}, \frac {x _ {2}}{x _ {0}}, \ldots , \frac {x _ {n}}{x _ {0}}\right)
\]

（E 到 \(\tilde{R}\) 中的映射）；由上一段所说的，这个函数可能不仅可以在 E 的闭包中的 \(R^{n}\) 的某些点处延拓，而且可以在 E 的闭包中 \(P_{n}\) 的某些“无穷远点”处延拓。

让我们证明，用这个方法可以得到例如代数中已经定义的一个实变量的有理函数到整个 \(\tilde{\mathbf{R}}\) 上的连续延拓。让我们把 \(\tilde{\mathbf{R}}\) 与 \(\mathbf{P}_1\) 等同起来，每个实数 \(x \in \mathbb{R}\) 等同于齐次坐标为 (1, x) 的点，而点 \(\infty\) 等同于齐次坐标为 (0,1) 的点。设 \(u(x)/v(x)\) 是一个有理函数，\(u\) 与 \(v\) 是次数分别为 \(m\) 与 \(n\) 的两个互素多项式；若例如假设 \(m \leqslant n\) ，并令 \(u_1(x,y) = x^n u(y/x)\) ，\(v_1(x,y) = x^n v(y/x)\) ，则有理函数 \(u/v\) ——在由实数 \(x\) （满足 \(v(x) \neq 0\) 的）组成的集上——可以看作映射 \((x,y) \to (v_1(x,y), u_1(x,y))\) 的限制。换句话说，我们把 \(u/v\) 连续延拓，给它值 \(\infty\) 于那些点 \(x \in \mathbb{R}\) ——使 \(v(x) = 0\) 的——处，并在点 \(\infty\) 处给它值 \(0\) （若 \(m < n\)）、值 \(\infty\) （若 \(m > n\)）、或首项系数之比（若 \(m = n\)）。

特别地，函数 \(1 / x\) 可以延拓到点 \(o\) ，在那里取值 \(\infty\) ，也可延拓到 \(\infty\) ，在那里取值 \(o\) ；这个延拓了的函数显然是 \(\tilde{\mathbf{R}}\) 到其自身的一个同胚。分式线性函数 \((ax + b) / (cx + d)\) （当 \(ad - bc \neq o\) 时）亦然。

同样，若 \(n\) 是 \(> 0\) 的整数，函数 \(x^n\) 可延拓到点 \(\infty\) ，在那里取值 \(\infty\)。

另一方面，两个实变量的有理函数一般不可能连续延拓到空间 \(\mathbf{P}_1 \times \mathbf{P}_1\) 或空间 \(\mathbf{P}_2\)（参见习题 5）。

